\let\negmedspace\undefined
\let\negthickspace\undefined
\documentclass[journal]{IEEEtran}
\usepackage[a5paper, margin=10mm, onecolumn]{geometry}
%\usepackage{lmodern} % Ensure lmodern is loaded for pdflatex
\usepackage{tfrupee} % Include tfrupee package

\setlength{\headheight}{1cm} % Set the height of the header box
\setlength{\headsep}{0mm}     % Set the distance between the header box and the top of the text

\usepackage{gvv-book}
\usepackage{gvv}
\usepackage{cite}
\usepackage{amsmath,amssymb,amsfonts,amsthm}
\usepackage{algorithmic}
\usepackage{graphicx}
\usepackage{textcomp}
\usepackage{xcolor}
\usepackage{txfonts}
\usepackage{listings}
\usepackage{enumitem}
\usepackage{mathtools}
\usepackage{gensymb}
\usepackage[breaklinks=true]{hyperref}
\usepackage{tkz-euclide} 
\usepackage{listings}
% \usepackage{gvv}                                        
\def\inputGnumericTable{}                                 
\usepackage[latin1]{inputenc}                                
\usepackage{color}                                            
\usepackage{array}                                            
\usepackage{longtable}                                       
\usepackage{calc}                                             
\usepackage{multirow}                                         
\usepackage{hhline}                                           
\usepackage{ifthen}                                           
\usepackage{lscape}
\usepackage{circuitikz}
\usepackage{comment}
\tikzstyle{block} = [rectangle, draw, fill=blue!20, 
    text width=4em, text centered, rounded corners, minimum height=3em]
\tikzstyle{sum} = [draw, fill=blue!10, circle, minimum size=1cm, node distance=1.5cm]
\tikzstyle{input} = [coordinate]
\tikzstyle{output} = [coordinate]


\begin{document}

\bibliographystyle{IEEEtran}
\vspace{3cm}

\title{12.144}
\author{EE25BTECH11026-Harsha}
 \maketitle
% \newpage
% \bigskip
{\let\newpage\relax\maketitle}

\renewcommand{\thefigure}{\theenumi}
\renewcommand{\thetable}{\theenumi}
\setlength{\intextsep}{10pt} % Space between text and floats


\numberwithin{equation}{enumi}
\numberwithin{figure}{enumi}
\renewcommand{\thetable}{\theenumi}

\textbf{Question}:\\
Let $\vec{A}=\myvec{3&&0&&0\\0&&6&&2\\0&&2&&6}$ and let $\lambda_1 \geq \lambda_2 \geq \lambda_3$ be the eigen values of $\vec{A}$.\\
\\
$\brak{a}$ The triple $\brak{\lambda_1,\lambda_2,\lambda_3}$ equals
\begin{enumerate}
\begin{multicols}{4}
    \item $\brak{9,4,2}$
    \item $\brak{8,4,3}$
    \item $\brak{9,3,3}$
    \item $\brak{7,5,3}$
\end{multicols}
\end{enumerate}
$\brak{b}$ The Matrix $\vec{P}$ such that
\begin{align*}
    \vec{P}^{-1}\vec{A}\vec{P}=\myvec{\lambda_1&&0&&0\\0&&\lambda_2&&0\\0&&0&&\lambda_3}
\end{align*}
is
\begin{enumerate}
\begin{multicols}{2}
    \item $\myvec{\frac{1}{\sqrt{3}}&&0&&\frac{-2}{\sqrt{6}}\\\frac{1}{\sqrt{3}}&&\frac{1}{\sqrt{2}}&&\frac{1}{\sqrt{6}}\\\frac{1}{\sqrt{3}}&&\frac{-1}{\sqrt{2}}&&\frac{1}{\sqrt{6}}}$
    \item $\myvec{\frac{1}{\sqrt{3}}&&\frac{-2}{\sqrt{6}}&&0\\\frac{1}{\sqrt{3}}&&\frac{1}{\sqrt{6}}&&\frac{1}{\sqrt{2}}\\\frac{1}{\sqrt{3}}&&\frac{1}{\sqrt{6}}&&\frac{-1}{\sqrt{2}}}$
    \item $\myvec{0&&0&&1\\\frac{1}{\sqrt{2}}&&\frac{1}{\sqrt{2}}&&0\\\frac{1}{\sqrt{2}}&&\frac{-1}{\sqrt{2}}&&0}$
    \item $\myvec{0&&1&&0\\\frac{1}{\sqrt{2}}&&0&&\frac{1}{\sqrt{2}}\\\frac{1}{\sqrt{2}}&&0&&\frac{-1}{\sqrt{2}}}$
\end{multicols}
\end{enumerate}
    
\solution \\
Let us solve the given question theoretically and then verify the solution computationally.\\
\\
$\brak{a}$ From the given matrix $\vec{A}$, as the value 3 is decoupled, i.e, the matrix $\vec{A}$ is invariant spanned under $\myvec{1&&0&&0}^{\top}$ ,
\begin{align}
    \therefore \lambda=3
\end{align}
Now the matrix reduces into $2 \times 2$ block,
\begin{align}
    \myvec{6&&2\\2&&6}
\end{align}
\begin{align}
    \myvec{6&&2\\2&&6}=4\myvec{1&&0\\0&&1}+2\myvec{1&&1\\1&&1}
\end{align}
\newpage
\vspace{0.25cm}
As we know that the eigen values of $\myvec{1&&1\\1&&1}$ is $\brak{2,0}$ and of $\myvec{1&&0\\0&&1}$ is $\brak{1,1}$, by definition of eigenvalues
\begin{align}
    \implies \myvec{6&&2\\2&&6}\vec{v}=\brak{4+2\mu}\vec{v}
\end{align}
where, $\vec{v}$ is the eigen-vector and $\mu$ is the eigen value(s) of $\myvec{1&&1\\1&&1}$ 
\begin{align}
    \therefore \lambda=\brak{4,8}
\end{align}
\begin{align}
    \implies \brak{\lambda_1,\lambda_2,\lambda_3}=\brak{8,4,3}
\end{align}

$\brak{b}$ The given relation can be computed as ,
\begin{align}
    \vec{A}=\vec{P}\vec{D}\vec{P}^{-1} \label{eq:1}
\end{align}
where $\vec{D}$ is the diagonal matrix of eigenvalues of $\vec{A}$.\\
\\
From ~\eqref{eq:1}, we can infer that it is the Eigen-value decomposition of matrix $\vec{A}$.\\
\\
Therefore, $\vec{P}$ is the ortho-normalized matrix of collection of eigen vectors of $\vec{A}$.
\begin{align}
    \vec{P}=\myvec{\vec{v_1}&&\vec{v_2}&&\vec{v_3}}
\end{align}
where $\vec{v_1},\vec{v_2} \,and\, \vec{v_3}$ are the normalized eigen vectors of $\vec{A}$.\\
From $\brak{a}$,
\begin{align}
    \vec{v_1}=\myvec{1\\0\\0} \qquad \vec{v_2}=\myvec{0\\\frac{1}{\sqrt{2}}\\\frac{1}{\sqrt{2}}} \qquad \vec{v_3}=\myvec{0\\\frac{1}{\sqrt{2}}\\\frac{-1}{\sqrt{2}}}
\end{align}
\begin{align}
    \vec{P}=\myvec{1&&0&&0\\0&&\frac{1}{\sqrt{2}}&&\frac{1}{\sqrt{2}}\\0&&\frac{1}{\sqrt{2}}&&\frac{-1}{\sqrt{2}}}
\end{align}

\end{document}

